\subsection{Higher-Order Tables and Sampling Design}

Two-way tables are the simplest categorical-data structure.  Higher-order tables introduce several possible notions of dependence: full independence, conditional independence given a third factor, marginal homogeneity, and interactions among factors.

\begin{definition} \label{def:two-way-log-linear-independence} \hlt{Two-Way Log-Linear Form}\\
Let $I,J\in\N$ and let $\pi_{ij}>0$ be two-way cell probabilities for $1\leq i\leq I$, $1\leq j\leq J$,\\
with $\pi_{i\cdot}\defeq\sum\limits_{j=1}^{J}\pi_{ij}$, $\pi_{\cdot j}\defeq\sum\limits_{i=1}^{I}\pi_{ij}$, $\alpha_i\defeq\log\pi_{i\cdot}$, and $\beta_j\defeq\log\pi_{\cdot j}$.
\begin{enumerate}[label=(\roman*)]
\setlength{\itemsep}{0pt}
\item \hlt{Independence form}: $\pi_{ij}=\pi_{i\cdot}\pi_{\cdot j}$ $\forall 1\leq i\leq I$, $1\leq j\leq J$ $\Longleftrightarrow$ $\log\pi_{ij}=\alpha_i+\beta_j$ $\forall 1\leq i\leq I$, $1\leq j\leq J$.
\item \hlt{Interaction terms}: $\gamma_{ij}\defeq\log\pi_{ij}-\alpha_i-\beta_j$, giving $\log\pi_{ij}=\alpha_i+\beta_j+\gamma_{ij}$.
\end{enumerate}
\hlt{Key Insight}: $\gamma_{ij}=0$ $\forall 1\leq i\leq I$, $1\leq j\leq J$ $\Longleftrightarrow$ $\pi_{ij}=\pi_{i\cdot}\pi_{\cdot j}$ $\forall 1\leq i\leq I$, $1\leq j\leq J$.
\end{definition}

\begin{remark} \label{rem:three-way-log-linear-models} \hlt{Three-Way Interactions}\\
Let $\pi_{ijk}>0$ for $(i,j,k)\in[I]\times[J]\times[K]$. A log-linear model with all two-factor interactions and no three-factor interaction is $\log\pi_{ijk}=\mu+\alpha_i+\beta_j+\eta_k+\delta_{ij}+\varepsilon_{ik}+\gamma_{jk}$.\\
The absent term $\lambda_{ijk}$ encodes the constraint of no third-order interaction; nested models are compared by deleting interaction terms and using the corresponding loss of degrees of freedom.
\end{remark}

\begin{principle} \label{prin:categorical-data-design-before-test} \hlt{Choose the Test from the Sampling Mechanism}\\
For a contingency table $N=(n_{ij})$, let $n_{i\cdot}$ and $n_{\cdot j}$ denote row and column margins.\\
Select the test from the design-fixed quantities:
\begin{enumerate}[label=(\roman*)]
\setlength{\itemsep}{0pt}
\item \hlt{Fixed Margins}: If $(n_{i\cdot})$ and $(n_{\cdot j})$ are fixed in a $2\times2$ table, use Fisher's exact test.
\item \hlt{Independent Samples}: If rows are independent multinomial samples with fixed $(n_{i\cdot})$,\\
test homogeneity of column probabilities.
\item \hlt{Single Multinomial Sample}: If $N\sim\operatorname{Multinomial}(n,(\pi_{ij}))$, test independence, i.e. $\pi_{ij}=\pi_{i\cdot}\pi_{\cdot j}$.
\item \hlt{Matched Pairs}: If observations are paired counts $n_{rs}$, test marginal homogeneity using discordant cells.
\item \hlt{Case-Control Sampling}: If $(n_{\cdot0},n_{\cdot1})$ are fixed by disease status, estimate odds ratios.\\
Absolute risks require population margins.
\end{enumerate}
\hlt{Key Insight}: The same numeric table can encode different null hypotheses under different sampling laws.
\end{principle}

\begin{remark} \label{rem:counts-not-percentages-in-chi-square-tests} \hlt{Use Counts, Not Percentages}\\
Let $Q(O,E)\defeq\sum_i\frac{(O_i-E_i)^2}{E_i}$. Pearson's chi-square calibration applies to counts: $Q(O/n,E/n)=\frac{1}{n}Q(O,E)$ and $Q(100O/n,100E/n)=\frac{100}{n}Q(O,E)$. Thus relative frequencies or percentages give a rescaled statistic unless the original count scale is restored.
\end{remark}
